% Standalone problem document, generated from the adjacent README.md.
% Compile with XeLaTeX. Mathematical content is maintained in Markdown.
\documentclass[11pt,a4paper]{article}
\usepackage[fontsize=10.5pt]{scrextend}
\usepackage[margin=22mm,headheight=15pt,headsep=9mm,footskip=11mm]{geometry}
\usepackage{fontspec}
\setmainfont{texgyrepagella}[Extension=.otf,UprightFont=*-regular,BoldFont=*-bold,ItalicFont=*-italic,BoldItalicFont=*-bolditalic]
\setsansfont{texgyreheros}[Extension=.otf,UprightFont=*-regular,BoldFont=*-bold,ItalicFont=*-italic,BoldItalicFont=*-bolditalic]
\setmonofont{texgyrecursor}[Extension=.otf,UprightFont=*-regular,BoldFont=*-bold,ItalicFont=*-italic,BoldItalicFont=*-bolditalic,Scale=MatchLowercase]
\usepackage{amsmath,amssymb}
\usepackage{unicode-math}
\setmathfont{texgyrepagella-math.otf}
\usepackage{xcolor,microtype,enumitem,needspace,fancyhdr}
\usepackage{bookmark}
\definecolor{ink}{HTML}{17344B}
\definecolor{accent}{HTML}{197D80}
\definecolor{muted}{HTML}{596773}
\hypersetup{colorlinks=true,linkcolor=accent,urlcolor=accent,pdftitle={IE-27 - Spectral
disk for a Radau stage
preconditioner},pdfauthor={Open Problems in Numerical Linear Algebra}}
\urlstyle{same}
\setlength{\parindent}{0pt}
\setlength{\parskip}{5pt plus 1pt minus 1pt}
\linespread{1.04}
\setlength{\emergencystretch}{3em}
\widowpenalty=10000
\clubpenalty=10000
\displaywidowpenalty=10000
\allowdisplaybreaks[1]
\setlist{nosep,leftmargin=1.5em}
\providecommand{\tightlist}{\setlength{\itemsep}{0pt}\setlength{\parskip}{0pt}}
\providecommand{\pandocbounded}[1]{#1}
\pagestyle{fancy}
\fancyhf{}
\fancyhead[L]{\sffamily\footnotesize\color{muted}OPEN PROBLEMS IN NUMERICAL LINEAR ALGEBRA}
\fancyhead[R]{\sffamily\bfseries\footnotesize\color{accent}IE-27}
\fancyfoot[L]{\parbox[b]{0.9\textwidth}{\sffamily\fontsize{8}{10}\selectfont\color{muted}Editorial ratings; a bounded search cannot certify that a problem remains unsolved.\\Literature check: 2026-09-11}}
\fancyfoot[R]{\sffamily\footnotesize\color{muted}\thepage}
\renewcommand{\headrulewidth}{0.3pt}
\renewcommand{\headrule}{\hbox to\headwidth{\color{accent}\leaders\hrule height \headrulewidth\hfill}}
\makeatletter
\renewcommand{\section}{\Needspace{5\baselineskip}\@startsection{section}{1}{0pt}{12pt plus 2pt minus 2pt}{4pt}{\sffamily\bfseries\large\color{ink}}}
\renewcommand{\subsection}{\Needspace{4\baselineskip}\@startsection{subsection}{2}{0pt}{9pt plus 1pt minus 1pt}{3pt}{\sffamily\bfseries\normalsize\color{ink}}}
\makeatother
\setcounter{secnumdepth}{0}
\begin{document}
{\sffamily\small\color{muted}Linear systems and elimination\par}
\vspace{3pt}
{\raggedright\hyphenpenalty=10000\exhyphenpenalty=10000\sffamily\bfseries\fontsize{21}{24}\selectfont\color{ink}Spectral
disk for a Radau stage preconditioner\par}
\vspace{7pt}
{\sffamily\small\textbf{Difficulty:} challenging\quad\textbf{Importance:} interesting
to specialist\par}
{\sffamily\small\bfseries\color{accent}Status: Partially resolved\par}
\vspace{4pt}
{\color{accent}\hrule height 0.8pt}
\vspace{6pt}
\textbf{Topic:} Preconditioning and eigenvalue localization

\textbf{Rating rationale:} The conjecture requires a sharp enclosure for
a nonnormal matrix family at every Runge--Kutta order and positive
shift. It would explain mesh- and time-step-independent spectral
clustering of a practical preconditioner.

\subsection{Statement}\label{statement}

Let \(q\ge2\), and let \(A_q\in\mathbb R^{q\times q}\) be the Butcher
matrix of the \(q\)-stage Radau IIA method. Explicitly, if
\(0<c_1<\cdots<c_q=1\) are the roots in \([0,1]\) of
\(P_q(2t-1)-P_{q-1}(2t-1)\), where \(P_j\) is the degree-\(j\) Legendre
polynomial normalized by \(P_j(1)=1\), and \(\ell_j\) are their Lagrange
cardinal polynomials, then

\[ (A_q)_{ij}=\int_0^{c_i}\ell_j(t)\,dt.\]

Take the exact factorization \(A_q^{-1}=L_qU_q\) without pivoting, with
\(L_q\) lower triangular and \(U_q\) upper triangular with unit
diagonal. Write \(\widehat U_q=U_q-I_q\), and assume
\(\|\widehat U_q\|_2<1\).

For \(n\ge1\), let \(M,K\in\mathbb R^{n\times n}\) be symmetric positive
definite mass and stiffness matrices from a finite-element
discretization of a coercive self-adjoint second-order elliptic
operator, and let \(\tau>0\). Define

\[\mathcal A=A_q^{-1}\otimes M+\tau I_q\otimes K,
\qquad\mathcal P_L=L_q\otimes M+\tau I_q\otimes K.\]

\textbf{Conjecture.} Every generalized eigenvalue \(\nu\in\mathbb C\) of
\(\mathcal Ax=\nu\mathcal P_Lx\) satisfies

\[|\nu-1|\le\|\widehat U_q\|_2.\]

The quantifiers are uniform over the stage number, admissible spatial
discretization and positive time step. This states the symmetric
elliptic setting used by the later spectral analysis explicitly. The
original Conjecture 1 says ``positive definite'' finite-element
matrices, and also discusses convection; this page does not silently
identify nonsymmetric positive definiteness with the symmetric setting.

\subsection{Equivalent small matrices and known
cases}\label{equivalent-small-matrices-and-known-cases}

The Kronecker structure reduces the spectrum to matrices

\[X_{q,\mu}=I_q+(I_q+\mu L_q^{-1})^{-1}\widehat U_q,
\qquad \mu\in\sigma(\tau M^{-1}K)\subset(0,\infty).\]

In particular, proving \(\rho(X_{q,\mu}-I_q)\le\|\widehat U_q\|_2\) for
every \(\mu>0\) proves the displayed conjecture for all admissible
discretizations. The source proves the two-stage case. At least \(n\)
eigenvalues equal one because \(\widehat U_q\otimes I_n\) has nullity at
least \(n\); this established part of the original conjecture is not an
additional target. Numerical spectra at larger stage counts support the
enclosure, but do not settle all \(q\) and \(\mu\).

\newpage
\subsection{References and status
check}\label{references-and-status-check}

\begin{itemize}
\tightlist
\item
  O. Axelsson, I. Dravins and M. Neytcheva, \emph{Stage-parallel
  preconditioners for implicit Runge--Kutta methods of arbitrarily high
  order, linear problems}, Numerical Linear Algebra with Applications
  31(1) (2024), e2532; first published 19 September 2023.
  \href{https://doi.org/10.1002/nla.2532}{DOI and full text};
  \href{https://www.diva-portal.org/smash/get/diva2\%3A1740890/FULLTEXT01.pdf}{institutional
  PDF}. §5.1, Conjecture 1, equations (17)--(18), and Theorem 1.
\item
  M. Outrata, \emph{On recent advances of spectral analysis for systems
  arising from fully-implicit RK methods}, Proceedings in Applied
  Mathematics and Mechanics (2026), e70082.
  \href{https://doi.org/10.1002/pamm.70082}{DOI};
  \href{https://arxiv.org/html/2510.21241v1}{author preprint
  arXiv:2510.21241v1}. §2.1, after equation (7), recalls the open
  general-stage conjecture; §2.3, equation (23), gives the small-matrix
  equivalence. The preprint says ``diameter'' at that point, whereas the
  original conjecture unambiguously specifies \textbf{radius}; the
  display follows the original.
\end{itemize}

On 2026-09-11, checked the original published conjecture and its
two-stage theorem, Outrata's 2025 preprint and 2026 publication record,
and searches for the title, Radau spectral-disk conjecture,
stage-parallel preconditioner proofs and counterexamples. No proof or
counterexample for arbitrary stage count was located. Later symbolic
spectral reductions are not a general proof of this particular disk
radius. The check is bounded. This is distinct from the
\href{https://github.com/ajt60gaibb/OpenProblemsInNLA/blob/main/linear-systems-and-elimination/IE-28/README.md}{diagonal
nilpotent-preconditioner existence question} and from the repository's
Krylov convergence and elimination-growth targets.
\end{document}
